\section{Background and Related Work}
\label{sec:ch3-related-work}
This section outlines architectural perspectives on experience, comfort studies in smart buildings, emerging studies of experience in HCI, and compares Affective Interaction (AfI) as an interpretive approach to affective computing. 

\subsection{Experiential Perspectives in Architecture}
\label{subsec:ch3-architecture-experience}
Architecture examines how built environments shape emotional life as relational, embodied, evolving, and co-constructed. Here, we review four influential strands: phenomenological accounts of embodied experience, environmental psychological perspectives, quantitative spatial traditions, and placemaking approaches.

\textbf{Phenomenological and embodied traditions} emphasise how architectural experience is multisensory and ``pre-reflective''. Here, Pallasmaa discusses how shadow, reverberation, texture, and warmth are felt before they are consciously interpreted~\cite{pallasmaa2024eyes}. Such experiences draw on embodied memories and ``primary images'' of shelter, movement, and materiality at an unconscious level. Zumthor frames architecture as an emotional encounter where light evokes bodily responses, invites calm, introspection, or awe~\cite{zumthor2006atmospheres}. Heidegger treats emotional experience as an existential being-in-the-world (\textit{dwelling}) situates shaped by place~\cite{heidegger1971building}. Scholarship on atmospheres highlights peripheral qualities like illumination, which impacts how spaces are interpreted and lived~\cite{zajonc1993catching, griffero2016atmospheres}. Collectively, these perspectives foreground diverse, embodied ways of experiencing and encourage occupants to form their own spatial narratives.

\textbf{Environmental psychology} examines how people emotionally appraise and interpret built environments. Classical work on appraisal highlights how spatial cues shape pleasure, arousal, and dominance~\cite{mehrabian1974approach}. Affordance research shows that spatial form invites or constrains action~\cite{vilar2012using}. Attention Restoration Theory and Stress Recovery Theory demonstrate how natural and architectural features impact stress, attentional recovery, and emotional regulation~\cite{basu2019attention, ulrich1991stress}. These perspectives show that affect arises from meaning, action possibilities, and anticipated effort within a space.

\textbf{Quantitative approaches} like space syntax~\cite{hillier2008space, hillier1976space} examine how spatial configuration shapes experience: integrated layouts foster lively environments, whereas segregated configurations evoke quieter, contemplative moods~\cite{jens2021design}. Here, research demonstrates that layout impacts individual usage patterns, social interactions, and thereby emotional tone.

Finally, \textbf{placemaking} supports emotional experience through cultural, social, and interpretive engagement with space. It bridges architectural and HCI perspectives by treating place as something shaped through interaction. Meaning is co-constructed through practice of using, personalising, and collectively inhabiting environments over time~\cite{dourish2001action, harrison1996re}. Emotional attachment, belonging, and atmosphere arise through these shared practices.


\subsection{Comfort Studies in Smart Buildings}
\label{subsec:ch3-comfort-studies}
Smart building research focuses on optimising comfort through environmental sensing~\cite{alavi2017comfort, alavi2020five}. This includes developing personalised thermal comfort via heating and cooling systems~\cite{jazizadeh2014human, khare2023thermal} as well as developing interactive systems such as indoor ``wind tunnels'' with small filaments that move to demonstrate presence of wind gusts~\cite{zhong2020hilo, snow2019performance, gaver2013indoor}. Studies on acoustic comfort typically gauge occupants' noise preferences through surveys~\cite{marques2020real, choi2009decision}, and visual comfort studies assess daylight availability and glare discomfort \cite{bian2023simulation, zheng2023daylighting}. These systems focus on adjusting indoor conditions, but treat users as passive recipients. Few studies have examined how comfort perceptions are shaped by social and cultural dynamics~\cite{clear2018thermokiosk}, shared control~\cite{von2021towards}, and behavioural patterns~\cite{clear2017d}. 

\subsection{Other Experiences in Smart Buildings}
Emerging works explore how smart technologies influence the ways people inhabit and interpret space~\cite{rao2025designing, wiberg2020interaction}. Displays manipulate light or sound to create calming, emotionally evocative environments~\cite{feng2019livenature, behrens2018sentiment, rao2023towards}. Speculative installations and design probes explore ambiguity, discovery, and future imaginaries (shared visions of possible futures)~\cite{gaver2006history, wiberg2020interaction}. For example, \emph{Pendaphonics}, a large-scale, physical–digital–sonic environment, used harmonic motion to invite playful, open-ended interaction and social improvisation among participants~\cite{hansen2009pendaphonics}. Similarly, the \emph{History Tablecloth} visualised accumulation of everyday actions in the home, inviting subjective reflection and social interaction~\cite{gaver2006history}.

Occupants’ bodies and behaviours are increasingly central. Studies have explored how movement and proximity affect interaction~\cite{homewood2021tracing, van2018skin}, and how systems can support connection and co-presence~\cite{wenhart2024there, hoare2015hide}. Others investigate neurodiverse design, showing how smart spaces can adapt to differing sensory needs~\cite{swaminathan2021tactile, garzotto2020interactive}. Collectively, they point to more experiential and socially situated understandings of smart environments. Yet, many still treat emotion in relatively instrumental or outcome-oriented terms. Our work builds upon this to foreground experience as a lived, interpretive, and relational process that is enacted through ongoing interaction with space, technology, and context.


\subsection{Affect and Affective Interaction (AfI)}
\label{sec:ch3-afi}
\paragraph{Affect}
In affective computing, affect is often understood as a measurable internal state, expressed through valence and arousal dimensions~\cite{russell1980circumplex} and captured via self-reports, physiology, or behavioural cues~\cite{picard1999affective}. Emotions in built environments have been studied using affective computing and neuroscience tools. These studies often rely on self-report, physiological sensors, or behavioural observations to infer affect~\cite{picard1999affective}. This framing supports smart building applications such as live emotion displays of occupants~\cite{behrens2018sentiment}, but assumes emotion is universal, discrete, and classifiable. However, affective science increasingly recognises emotion as ambiguous, contextual, and socially constructed~\cite{feldman2022affect, barrett2022context}. This shift motivates more context-sensitive approaches that treat affect as a relational and meaning-laden phenomenon. Our work responds to this shift by adopting an interpretivist stance.

\paragraph{Affective Interaction (AfI)}
AfI offers an alternative to state-based affect models. Rather than treating affect as input/output, AfI foregrounds the interactional unfolding of affect over time~\cite{lottridge2011affective}. Affect is shaped through dialogue between user and system~\cite{hook2018designing, hook2009affective}, interpretive flexibility~\cite{boehner2005affect, boehner2007emotion}, and personal reflection~\cite{ahmadpour2025affective}. While AfI has been explored in domestic~\cite{angelini2015towards}, urban~\cite{paraschivoiu2022affective},  creative~\cite{huang2025affective}, and therapeutic contexts~\cite{ahmadpour2025affectivevr}, it remains underutilised in smart building design~\cite{rao2025lessons}. Buildings are affective spaces: occupants reflect on their bodies and surroundings, and co-construct meaning through daily use. By applying AfI in HBI, we aim to surface these everyday affective practices and rethink how smart environments might support---not override---emotionally meaningful experience.
